%% =========================================================
%%  Conformance-Gated Generation for Enterprise Code Synthesis
%%  Journal of Systems and Software (JSS), Elsevier
%%  Manuscript Type: Research Article (Empirical Study)
%% =========================================================
\documentclass[review,12pt,3p]{elsarticle}

%% ----- Packages ------------------------------------------
\usepackage{float}
\usepackage{amsmath,amssymb}
\usepackage{booktabs}
\usepackage{multirow}
\usepackage{array}
\usepackage{graphicx}
\usepackage{xcolor}
\usepackage{hyperref}
\usepackage{lmodern}
\usepackage{microtype}
\usepackage{setspace}
\usepackage{lineno}
\usepackage{url}
\usepackage[capitalise,noabbrev]{cleveref}
\usepackage{listings}
\usepackage{tcolorbox}
\tcbuselibrary{skins,breakable}
\usepackage[table]{xcolor}

\lstdefinelanguage{JavaScript}{
  keywords={break, case, catch, continue, debugger, default, delete, do, else, finally, for, function, if, in, instanceof, new, return, switch, this, throw, try, typeof, var, void, while, with, let, const, export, import},
  morecomment=[l]{//},
  morecomment=[s]{/*}{*/},
  morestring=[b]',
  morestring=[b]",
  sensitive=true
}

\doublespacing
\modulolinenumbers[5]

%% ----- Journal metadata ----------------------------------
\journal{Journal of Systems and Software}

%% =========================================================
\begin{document}

\begin{frontmatter}

\title{Conformance-Gated Generation for Enterprise Code Synthesis:\\
An Empirical Study of Model-Dependent Effects Across LLM Families}

\author[inst1]{Raduanul Haque Maruf}
\ead{2023331001@student.sust.edu}

\author[inst2]{Pabel Shahrear\corref{cor1}}
\ead{shahrear-mat@sust.edu}

\cortext[cor1]{Corresponding author}

\affiliation[inst1]{organization={Department of Computer Science and Engineering},
    institution={Shahjalal University of Science and Technology},
    city={Sylhet},
    postcode={3114},
    country={Bangladesh}}

\affiliation[inst2]{organization={Department of Mathematics},
    institution={Shahjalal University of Science and Technology},
    city={Sylhet},
    postcode={3114},
    country={Bangladesh}}

%% --------- Abstract --------------------------------------
\begin{abstract}
Large language models (LLMs) increasingly serve as code generation engines within
enterprise software pipelines. Unrestricted generation produces architecturally unsafe,
verbose outputs that violate project-specific structural invariants, a phenomenon we
term \emph{architectural drift}. This paper presents the Aether Framework, a
conformance-gated generation system that injects a machine-readable architectural
contract into the LLM prompt and evaluates generated code against a custom rule engine
(the G-Suite). We report a study of 2,451 evaluations across 30 enterprise backend
prompts, two conditions (Vanilla versus Aether), and five model families ranging from
7B to over 100B parameters. The central finding is a capacity-dependent funnel effect
with a threshold between 8B and 12B parameters. Below this threshold, Qwen 2.5 Coder
7B shows a small negative effect ($d = -0.36$) and Llama 3.1 8B shows an effect
indistinguishable from zero ($d = -0.010$). Above it, Gemma 3 12B ($d = +1.62$),
Qwen 2.5 Coder 14B ($d = +1.33$), and Gemini 2.5 Flash ($d = +1.26$) show large,
significant reductions in output tokens. Structural correctness is maintained at
95.1\% (Gemma 3 12B Aether: 66.3\%, Llama 3.1 8B Aether: 93.3\%, all others: 100\%),
though Aether's compression coincides with reduced compilation reliability for
mid-capacity models, revealing a reliability trade-off. Generation entropy, via
coefficient of variation, paradoxically increases under Aether for three of four
local open-weight models. These results provide the first empirical evidence that
architectural pre-conditioning exhibits capacity-dependent effects, acting as an
effective output funnel above 12B parameters while ambiguating smaller quantized
models. Practical implications for enterprise LLM deployment are discussed.
\end{abstract}

\begin{keyword}
LLM code generation \sep
architectural conformance \sep
prompt engineering \sep
token efficiency \sep
enterprise software \sep
empirical study \sep
model-dependent effects
\end{keyword}

\end{frontmatter}

\linenumbers

%% =========================================================
\section{Introduction}
\label{sec:intro}
%% =========================================================

The adoption of large language models as practical code generation engines has accelerated
substantially in recent years, with tools such as GitHub Copilot~\citep{Copilot2022},
Amazon CodeWhisperer~\citep{CodeWhisperer2023}, and Google Gemini Code Assist demonstrating
tangible productivity gains across a range of programming tasks~\citep{Ziegler2022Productivity}.
However, as enterprise organisations integrate these systems into production pipelines, a
critical and underexplored class of failure has emerged that goes beyond syntactic
correctness: \emph{architectural drift}. LLMs, operating without project-specific
structural context, exhibit a characteristic tendency to generate self-contained,
boilerplate-heavy code that assumes a generic project environment. This behaviour produces
outputs that violate project-specific constraints such as module boundary isolation,
mandatory authentication guard placement, and database driver scoping to designated
repository layers. The consequence is not merely aesthetic --- it is structural, financial,
and operational. Architecturally drifted code must be manually reviewed and refactored
before it can be safely integrated, while the verbosity of unconstrained generation
translates directly into increased token consumption and therefore elevated API inference
costs.

Prior empirical work on LLM code generation has focused predominantly on functional
correctness, evaluating models against unit test suites~\citep{Chen2021HumanEval,
Austin2021MBPP} or canonical competitive programming
benchmarks~\citep{Li2022AlphaCode}. Studies examining prompt engineering strategies have
demonstrated that structured approaches --- including chain-of-thought
prompting~\citep{Wei2022ChainOfThought}, few-shot exemplar injection~\citep{Brown2020GPT3},
and system-level context prefixes --- can meaningfully improve the quality of generated
outputs. However, to our knowledge, no prior study has systematically measured the effect
of injecting a \emph{strict, machine-readable architectural contract} into the generation
prompt on output token efficiency, API cost, and generation consistency across five architecturally
distinct model families spanning from 7B to 100B+ parameters. This gap represents a significant
blind spot in the empirical literature, particularly as enterprise organisations seek to deploy
LLM code generation at scale under cost and safety constraints.

This paper addresses that gap through the design, implementation, and empirical evaluation
of the \emph{Aether Framework} --- a conformance-gated generation system built specifically
for enterprise backend code synthesis. Aether operates on two primary design principles.
First, it enforces a \emph{context-first} generation discipline: before any code is
generated, a structured architectural contract is assembled and injected into the prompt.
This contract encodes module boundaries, permissible import paths, authentication
requirements, and generation constraints in a machine-readable format. Second, Aether
implements a \emph{conformance gate}: generated outputs are evaluated against a custom
rule engine (the G-Suite) that maps each structural invariant to a published architectural
standard, and non-conforming outputs are subjected to a self-healing regeneration loop.

The study is organised around four research questions:

\begin{description}
  \item[\textbf{RQ1}] Does architectural pre-conditioning via the Aether Framework affect
    structural code correctness compared to standard prompting?
  \item[\textbf{RQ2}] Does architectural pre-conditioning significantly reduce LLM output
    token generation, and does this effect differ across model families?
  \item[\textbf{RQ3}] Does architectural pre-conditioning reduce weighted API cost,
    accounting for both input context overhead and output generation volume?
  \item[\textbf{RQ4}] Does architectural pre-conditioning improve generation consistency
    across repeated runs on the same prompt?
\end{description}

The contributions of this work are fourfold. First, we introduce and describe the Aether
Framework, including the \texttt{module.contract.ts} and \texttt{ai-context.md} artefacts
and the G-Suite rule engine. Second, we release a 2,451-evaluation empirical dataset across
five model families and two conditions. Third, we report the first empirical evidence that
architectural pre-conditioning produces \emph{model-dependent effects} --- significantly
beneficial for large closed-weight models and neutral or harmful for small open-weight
quantized models. Fourth, we document the P30 qualitative finding as a concrete
demonstration of catastrophic context overflow prevention. The remainder of this paper is
structured as follows: \cref{sec:background} reviews related work; \cref{sec:framework}
describes the Aether Framework; \cref{sec:methodology} details the experimental methodology;
\cref{sec:results} presents and interprets the results; \cref{sec:discussion} discusses
implications; and \cref{sec:conclusion} concludes with directions for future work.

%% =========================================================
\section{Background and Related Work}
\label{sec:background}
%% =========================================================

\subsection{LLM-Based Code Generation}

The emergence of transformer-based models pre-trained on large corpora of source code has
fundamentally altered the landscape of automated program synthesis. Chen et
al.~\citep{Chen2021HumanEval} introduced HumanEval, a functional correctness benchmark
comprising 164 handwritten programming problems, and demonstrated that the Codex model
(12B parameters) achieved a pass@1 rate of 28.8\%, establishing a foundational metric for
the field. Subsequent work by Austin et al.~\citep{Austin2021MBPP} contributed the MBPP
dataset and showed that larger models exhibit substantially improved synthesis capabilities
on diverse programming tasks. Li et al.~\citep{Li2022AlphaCode} achieved competitive
performance on Codeforces programming contests through a specialised training regime
combining competitive programming data and systems-level sampling strategies. These studies
share a common evaluation paradigm: \emph{functional correctness} is measured by executing
generated code against unit test suites or comparing outputs against reference solutions.
Structural properties of the generated code --- its compliance with project-level
architectural patterns, its cost as measured in inference tokens, and its consistency across
repeated invocations --- receive markedly less attention.

\subsection{Prompt Engineering and Context Injection}

A substantial body of work has examined how the structure of the input prompt affects LLM
output quality. Wei et al.~\citep{Wei2022ChainOfThought} demonstrated that chain-of-thought
prompting, which encourages intermediate reasoning steps, significantly improves performance
on complex reasoning tasks. Brown et al.~\citep{Brown2020GPT3} established the utility of
few-shot exemplar injection, showing that models can rapidly adapt to novel tasks when
provided with a small number of task-specific examples in the prompt. Liu et
al.~\citep{Liu2023PromptSurvey} provide a comprehensive survey of prompt engineering
strategies, documenting the breadth of techniques applied to code generation, question
answering, and instruction following. More recently, structured prompting approaches have
explored the systematic encoding of role assignments, persona instructions, and task
decomposition within the prompt~\citep{Jiang2020StructuredPrompt}. However, these
approaches are generally evaluated on functional correctness benchmarks and do not address
the injection of \emph{architectural} structural constraints, nor do they examine the
differential effect of structured prompting across models of substantially different
capacities.

\subsection{Software Architecture Conformance Checking}

The problem of ensuring that implemented code conforms to an intended architecture has a
long history in software engineering research. Traditional approaches rely on static
analysis tools that compare the implemented module dependency graph against an architectural
specification~\citep{deGraft2009ArchConformance}. Tools such as ArchUnit~\citep{ArchUnit2021}
and similar constraint-checking frameworks operationalise architectural rules as executable
test cases that can be integrated into a continuous integration pipeline. These approaches
operate exclusively on existing code; they are not designed to influence the generation
process of an AI system. The Aether Framework bridges this gap by encoding architectural
constraints as prompt-level context, influencing generation \emph{prospectively} rather
than detecting violations \emph{retrospectively}.

\subsection{Token Efficiency in LLM Inference}

The computational and financial cost of LLM inference is directly proportional to the
number of tokens processed and generated~\citep{Sharir2020CostNLP}. As organisations
deploy code generation at scale, the token efficiency of generation strategies becomes
a significant operational concern. Verbose generation increases billing costs on commercial
API providers and increases latency, negatively affecting developer experience in interactive
settings. Despite its practical importance, token efficiency has received limited attention
as an independent evaluation dimension in empirical code generation studies. Work on
speculative decoding~\citep{Leviathan2023Speculative} and model quantization~\citep{Dettmers2022LLMInt8}
addresses efficiency from the model architecture and serving stack perspective, but does not
examine how prompt-level strategies can reduce the volume of generated tokens. This paper
directly addresses this gap.

\subsection{Gap Statement}

The reviewed literature reveals a consistent pattern: empirical evaluations of LLM code
generation prioritise functional correctness, while structural safety, token efficiency, and
generation consistency remain largely unmeasured. Prompt engineering research demonstrates
that structured context injection can improve output quality but does not examine
\emph{architectural} constraints or their differential effect across model families. No
prior work, to our knowledge, has measured the effect of injecting strict, machine-readable
architectural contracts on output token volume, weighted API cost, and cross-iteration
generation entropy, nor has any study reported whether such effects are consistent across
large closed-weight and small open-weight quantized models. This paper fills that gap.

%% =========================================================
\section{The Aether Framework}
\label{sec:framework}
%% =========================================================

The Aether Framework is an AI-first web framework designed to enforce architectural
conformance during LLM code generation. Its central mechanism --- the
\emph{conformance-gated generation loop} --- departs from the standard code generation
paradigm by requiring that architectural context be fully assembled and injected into the
prompt \emph{before} the model generates any code. This section describes the three
components of the framework: the contract artefacts, the G-Suite rule engine, and the
generation loop itself.

\subsection{Architectural Contract Artefacts}

The Aether contract consists of two machine-readable files that together encode the
structural invariants of the target module.

The \texttt{module.contract.ts} file is a TypeScript module that programmatically declares
the boundaries, dependencies, and exposed surface of a single software module using a
\texttt{defineContract()} API. Concretely, it enumerates the routes exposed by the module
(including HTTP method, path, and authentication requirements), the functions available to
other modules, the types exported by the module, the declared side effects (such as database
table reads and writes), and the invariants that must hold for every generated implementation.
Crucially, the contract also specifies the module's dependencies on other modules in the
system, forming the basis for the cross-module import rules enforced by the rule engine.


Listing~\ref{lst:contract} shows an excerpt from the \texttt{auth} module contract used
in this study.

\begin{lstlisting}[language=JavaScript, caption={Excerpt from \texttt{modules/auth/module.contract.ts}}, label=lst:contract, basicstyle=\small\ttfamily, frame=single, breaklines=true]
export default defineContract({
  name: 'auth',
  exposes: {
    routes: [
      { method: 'POST', path: '/auth/login',   auth: false },
      { method: 'POST', path: '/auth/logout',  auth: true  },
    ],
    functions: ['requireUser', 'createSession', 'invalidateSession'],
  },
  invariants: ['Expired sessions must never be returned as valid'],
});
\end{lstlisting}

The \texttt{ai-context.md} file is a Markdown document that provides natural-language and
structured architectural rules to the LLM. It describes the project's technology stack,
the permitted database access patterns (i.e.\ the repository layer abstraction), the
authentication requirement for protected routes, and the expected output format for
generated files. Critically, the \texttt{ai-context.md} instructs the model to generate
\emph{only} the code artefacts required to implement the requested feature within the
module's declared boundaries, suppressing the model's tendency to generate surrounding
boilerplate, mock data layers, and standalone server bootstrap code.

The Vanilla baseline condition uses an identical prompt structure and the same feature
request but \emph{omits} the contract injection entirely. The model receives only the
feature description and edge case specification, receiving no information about module
boundaries, permitted imports, authentication requirements, or expected output scope.

\subsection{The G-Suite Rule Engine}

The G-Suite is a custom evaluation engine that assesses generated code against five
architectural rules. Each rule is mapped to a canonical published standard to ensure
construct validity.

Rule G1 prohibits direct database driver imports outside the designated \texttt{lib/db.ts}
layer, operationalising the Repository Pattern as described by Fowler~\citep{Fowler2002PEAA}.
This rule ensures that routing layer code interacts with data exclusively through the
repository abstraction rather than through raw database client calls.

Rule G2 verifies that all routes declared as authenticated in the module contract include
an explicit call to the mandatory authentication guard function (\texttt{requireUser}),
operationalising the secure-by-default principle from the OWASP Secure Coding Practices
guidelines~\citep{OWASP2023Secure}.

Rule G3 verifies that no module imports any function, type, or constant from a module not
listed in its declared dependencies, operationalising the Dependency Rule of Clean
Architecture as formulated by Martin~\citep{Martin2017CleanArch}.

Rule G4 is reserved for invariant test execution and was deferred to a subsequent phase of
the project; consequently, it is skipped in all evaluations reported here and excluded from
pass-rate calculations.

Rule G5 verifies that the generated TypeScript code compiles without type errors under the
project's \texttt{tsconfig.json}, ensuring basic syntactic and type-level correctness.

\subsection{The Generation Loop}

The generation loop proceeds as follows. Given a feature request, the framework assembles
the prompt by concatenating the project manifest, the target module's
\texttt{module.contract.ts} content, the \texttt{ai-context.md} content, the feature
description, and the edge case specification. The assembled prompt is submitted to the
configured LLM provider. The response is parsed to extract generated source files, which
are then evaluated by the G-Suite rule engine. If all non-skipped rules pass, the loop
terminates and the generated files are recorded as the output of the evaluation. If any
rule fails, an error hint summarising the failing rules and their remediation guidance is
assembled and injected into the subsequent attempt's prompt, up to a maximum of three
attempts. The Vanilla baseline executes the same loop but without the contract injection
step and with the contract suite check disabled.

%% INSERT Figure 1: Aether Framework Conformance-Gated Loop
\begin{figure}[H]
  \centering
  % INSERT fig1_conformance_loop below (export from architecture_diagrams.md Figure 1)
   \includegraphics[width=\linewidth]{figures/fig1_conformance_loop.pdf}
  \caption{The Aether conformance-gated generation loop. A feature request is processed
  under two parallel paths: the Aether path assembles the architectural contract
  (\texttt{module.contract.ts} + \texttt{ai-context.md}) and injects it into the prompt,
  while the Vanilla baseline path omits the contract. In both cases, the G-Suite rule
  engine evaluates the generated output; under Aether, rule failures trigger a
  self-healing retry loop (up to three attempts) with an error hint appended to the prompt.}
  \label{fig:loop}
\end{figure}

%% =========================================================
\section{Methodology}
\label{sec:methodology}
%% =========================================================

\subsection{Dataset Design}

The evaluation dataset consists of 30 enterprise backend feature prompts distributed across
seven functional domains: authentication (6 prompts), task management (9 prompts), billing
(3 prompts), team management (4 prompts), notifications (3 prompts), reports (2 prompts),
and file management (3 prompts). Each prompt specifies a concrete HTTP route to be
implemented (e.g.\ \texttt{POST /auth/register}, \texttt{GET /tasks}) along with a
description of the required behaviour and at least one edge case that the generated code
must handle. Prompts are assigned a complexity level --- low (8 prompts), medium (16
prompts), or high (6 prompts) --- based on the number of distinct operations, the number
of interacting modules, and the specificity of the edge case.

The domain and complexity distribution was designed to reflect the structure of a realistic
enterprise SaaS backend. Authentication and task management prompts constitute the majority
of the dataset, consistent with their central role in business application architectures.
The inclusion of file management, billing, and notification domains ensures that the
evaluation covers cross-cutting concerns such as third-party service integration and event
emission.
%% INSERT Table 1: Benchmark Dataset (publication_tables.tex → Table 1)
% --- Table 1: Benchmark Dataset ---
\begin{table*}[htbp]
\centering
\small
\setlength{\tabcolsep}{4pt}
\caption{Benchmark Dataset consisting of 30 enterprise backend routing scenarios across multiple business domains and complexity levels. Descriptions abbreviated; full text in
replication package (Data Availability).}
\label{tab:dataset}
\resizebox{\linewidth}{!}{
\begin{tabular}{llllll}
\toprule
\textbf{Prompt ID} & \textbf{Domain} & \textbf{Route} & \textbf{HTTP Method} & \textbf{Complexity} & \textbf{Description} \\
\midrule
P01 & Auth & \texttt{/auth/register} & POST & Low & Create a user registration route that validates email, hashes the password, a... \\
P02 & Auth & \texttt{/auth/login} & POST & Low & Create a login route that takes email and password, looks up the user, and re... \\
P03 & Auth & \texttt{/auth/logout} & POST & Low & Create a logout route that invalidates the current session and clears the coo... \\
P04 & Tasks & \texttt{/tasks} & GET & Low & Create a list-tasks route that returns all tasks for the current user's team,... \\
P05 & Tasks & \texttt{/tasks} & POST & Medium & Create a create-task route that accepts { title, description, assigneeId }, v... \\
P06 & Tasks & \texttt{/tasks/\:id/status} & PATCH & Medium & Create a status-update route that transitions a task between 'todo', 'in\_prog... \\
P07 & Tasks & \texttt{/tasks/\:id/assign} & POST & Medium & Create an assign-task route that takes assigneeId, verifies the assignee is a... \\
P08 & Tasks & \texttt{/tasks/\:id} & DELETE & Medium & Create a soft-delete route for tasks. The task must remain in the database wi... \\
P09 & Tasks & \texttt{/tasks/\:id/archive} & POST & Medium & Create an archive route that moves finished tasks older than 30 days to an ar... \\
P10 & Tasks & \texttt{/tasks/\:id/history} & GET & Medium & Create a route that returns the audit-log history of a single task (status ch... \\
P11 & Tasks & \texttt{/tasks/\:id/transfer-team} & POST & High & Create a route that transfers a task from one team to another. Both teams mus... \\
P12 & Tasks & \texttt{/tasks/bulk} & PATCH & High & Create a bulk-update route that accepts an array of {id, patch} objects, vali... \\
P13 & Billing & \texttt{/billing/subscriptions} & POST & Medium & Create a route that creates a Stripe subscription for the current user, stori... \\
P14 & Billing & \texttt{/billing/summary} & GET & Medium & Create a loader route that returns a billing summary for the current user: cu... \\
P15 & Billing & \texttt{/billing/webhook} & POST & High & Create a Stripe webhook handler that verifies the signature, parses the event... \\
P16 & Auth & \texttt{/auth/refresh} & POST & Medium & Create a token-refresh route that exchanges a refresh token for a new access ... \\
P17 & Auth & \texttt{/auth/password-reset/request} & POST & Medium & Create a password-reset-request route that emails a one-time reset link. Alwa... \\
P18 & Auth & \texttt{/auth/password-reset/confirm} & POST & Medium & Create a password-reset-confirm route that consumes a one-time token, sets th... \\
P19 & Teams & \texttt{/teams} & POST & Medium & Create a route that creates a new team, with the caller as the owner. Owner i... \\
P20 & Teams & \texttt{/teams/\:id/invite} & POST & Medium & Create an invite-member route that generates a single-use invite token, email... \\
P21 & Teams & \texttt{/teams/\:id/members/\:userId} & DELETE & Medium & Create a remove-member route. The owner cannot be removed. After removal, all... \\
P22 & Teams & \texttt{/teams/\:id} & PATCH & Low & Create a route to rename a team. Only admins may call it. \\
P23 & Notifications & \texttt{/notifications/send} & POST & Medium & Create a route that sends a notification to a list of user ids. The notificat... \\
P24 & Notifications & \texttt{/notifications} & GET & Low & Create a route that lists the current user's notifications, paginated, with u... \\
P25 & Notifications & \texttt{/notifications/\:id/read} & POST & Low & Create a mark-as-read route. Idempotent — calling twice must not error. \\
P26 & Reports & \texttt{/reports/team-velocity} & GET & High & Create a route that returns the team's task velocity over the last 30 days: t... \\
P27 & Reports & \texttt{/reports/billing-forecast} & GET & High & Create a route that returns a 3-month billing forecast based on current subsc... \\
P28 & Files & \texttt{/files/upload} & POST & High & Create a file-upload route that accepts a multipart form, validates size and ... \\
P29 & Files & \texttt{/files/\:id} & GET & Medium & Create a download route that streams a file from object storage. Caller must ... \\
P30 & Files & \texttt{/files/\:id} & DELETE & Medium & Create a soft-delete route for files. Removes the row, deletes the underlying... \\
\bottomrule
\end{tabular}
}
\end{table*}
\subsection{Model Selection}

Five models were selected to represent distinct points across the LLM capability spectrum, ranging from 7B to over 100B parameters. Gemini 2.5 Flash is a large, closed-weight model offered by Google via a commercial API, representative of the managed services enterprise organisations typically consume. The remaining four models are open-weight and were deployed locally via the Ollama inference server with 4-bit K-matrix quantization (Q4\_K\_M). These comprise: Qwen 2.5 Coder 7B, Llama 3.1 8B, Gemma 3 12B, and Qwen 2.5 Coder 14B.

This set of five models was chosen deliberately to stress-test the generalisability of the Aether approach across the capacity divide. They differ not only in parameter count but also in training data composition, context window capacity, and inference characteristics. If the Aether effect were uniform across this spectrum, it would provide strong evidence of a model-agnostic mechanism. By evaluating a gradient of parameter scales, this comparison aims to reveal the precise capacity threshold at which architectural pre-conditioning becomes beneficial.

%% INSERT Table 2: Model Specifications (publication_tables.tex → Table 2)
% --- Table 2: Model Specifications ---
\begin{table*}[htbp]
\centering
\caption{Specifications and execution environment characteristics of the evaluated large language models (LLMs), ordered by parameter scale descending.}
\label{tab:models}
\resizebox{\linewidth}{!}{
\begin{tabular}{lllllll}
\toprule
\textbf{Model} & \textbf{Parameters} & \textbf{Type} & \textbf{Quantization} & \textbf{Deployment} & \textbf{Repetitions ($n$)} & \textbf{Context Window} \\
\midrule
Gemini 2.5 Flash & \(\sim\)100B+ & Closed-weight & Full precision & Google API & 1 & 1M tokens \\
Qwen 2.5 Coder 14B & 14B & Open-weight & Q4\_K\_M & Local (Ollama) & 10 & 8192 tokens \\
Gemma 3 12B & 12B & Open-weight & Q4\_K\_M & Local (Ollama) & 10 & 8192 tokens \\
Llama 3.1 8B & 8B & Open-weight & Q4\_K\_M & Local (Ollama) & 10 & 8192 tokens \\
Qwen 2.5 Coder 7B & 7B & Open-weight & Q4\_K\_M & Local (Ollama) & 10 & 8192 tokens \\
\bottomrule
\end{tabular}
}
\end{table*}

\subsection{Experimental Infrastructure}

The evaluation pipeline was implemented entirely in TypeScript and executed in a headless,
fully automated mode. The pipeline incorporates exponential backoff with jitter for API
rate limit recovery and per-evaluation state checkpointing so that interrupted runs can be
resumed without loss of data. Each completed evaluation is serialised to a JSON file
recording the full prompt, all attempt metrics (tokens in, tokens out, latency, rule
outcomes, and generated files), and the final pass status.

Gemini 2.5 Flash was accessed via the Google Generative AI API over a standard internet
connection. Qwen 2.5 Coder 7B Q4\_K\_M was served by an Ollama instance running on an
NVIDIA RTX 3050 8 GB GPU workstation, accessed from the evaluation laptop via a Tailscale
VPN tunnel. This configuration introduces a fixed network overhead to Qwen latency
measurements; all latency figures should therefore be interpreted as \emph{wall-clock}
values inclusive of VPN transit and Ollama scheduling, not as pure model inference time.

The evaluation pipeline was developed with the assistance of an AI coding agent (Google
Antigravity), which supported implementation of the TypeScript codebase. All experimental
design decisions, architectural logic, and evaluation methodology were directed and
validated by the authors.

\subsection{Experimental Procedure}

For Gemini 2.5 Flash, each of the 30 prompts was evaluated once under each condition,
yielding 60 evaluations (30 Vanilla + 30 Aether). The single-repetition design for the
commercial API was necessitated by cost constraints.

For the four local open-weight models (Qwen 7B, Llama 8B, Gemma 12B, and Qwen 14B), each
of the 30 prompts was evaluated ten times under each condition, yielding a nominal total
of 600 evaluations per model. The ten-repetition design serves a dual purpose: it enables
the estimation of generation entropy (within-prompt variation) and provides a more stable
estimate of the mean output token count for the Wilcoxon signed-rank test, which is conducted
on prompt-level means.

Data exclusion logic was applied to two specific anomalies to preserve the integrity of the
paired statistical tests. First, for Qwen 2.5 Coder 7B, prompt P30 under the Vanilla condition
resulted in a degenerate generation: the first Vanilla iteration immediately exhausted the
local 4,096-token context window, halting further data collection after a single overflow
response. Because the Wilcoxon test requires paired observations, the entire P30
prompt (both Vanilla and Aether) was dropped for Qwen 7B, leaving 29 valid prompt pairs.
One iteration of prompt P05 under the Aether condition for Llama 3.1 8B produced output
truncated at the 8,192-token context limit and was excluded from aggregate computation.
The remaining nine iterations for this prompt were retained, preserving the P05 prompt pair
in the Wilcoxon analysis (N=30 for Llama 3.1 8B). Across all five models, the study
comprises a total of 2,451 valid individual evaluations.

%% INSERT Figure 2: End-to-End Experimental Pipeline
\begin{figure}[H]
  \centering
  % INSERT fig2_pipeline below (export from architecture_diagrams.md Figure 2)
   \includegraphics[width=\linewidth]{figures/fig2_pipeline.pdf}
  \caption{End-to-end experimental pipeline. Thirty prompts spanning seven domains and
  three complexity levels are forked into Vanilla and Aether conditions. Each of the five
  evaluated LLMs processes both conditions (n=1 per prompt for Gemini 2.5 Flash; n=10 per
  prompt for the four local models). Generated outputs pass through the G-Suite rule engine
  for metrics collection, followed by statistical analysis (Wilcoxon signed-rank test,
  Cohen's $d$, coefficient of variation).}
  \label{fig:pipeline}
\end{figure}

\subsection{Metrics}

Five primary metrics were collected for each evaluation. \emph{Tokens in}
(\texttt{tokens\_in}) records the number of tokens in the assembled prompt as reported by
the respective model API or inference server. \emph{Tokens out} (\texttt{tokens\_out})
records the number of tokens in the generated response. \emph{Latency} (\texttt{latency\_ms})
records the wall-clock time from prompt submission to response receipt in milliseconds.
\emph{Weighted cost} is a derived metric defined in \cref{eq:wcost}. \emph{Structural
pass rate} is the proportion of evaluations in which all non-skipped G-Suite rules pass.
\emph{Generation entropy}, defined in \cref{eq:entropy}, is computed for all four local
open-weight models (Qwen 2.5 Coder 7B, Llama 3.1 8B, Gemma 3 12B, and Qwen 2.5 Coder 14B),
where multiple iterations per prompt exist.

The weighted cost metric is defined as:
\begin{equation}
  \text{weighted\_cost} = (\text{tokens\_in} \times 1) + (\text{tokens\_out} \times 4)
  \label{eq:wcost}
\end{equation}

The factor of four applied to output tokens reflects the approximate pricing ratio between
input and output tokens on commercial LLM API providers as of the time of writing, where
output token generation is computationally more expensive than input token processing.

Generation entropy is operationalised as the coefficient of variation of \texttt{tokens\_out}
across iterations for a given prompt under a given condition:
\begin{equation}
  \text{entropy}_{p,c} = \frac{\sigma(\text{tokens\_out}_{p,c})}{\mu(\text{tokens\_out}_{p,c})}
  \label{eq:entropy}
\end{equation}
where $\sigma$ denotes the standard deviation and $\mu$ denotes the mean across iterations
for prompt $p$ under condition $c$. The mean entropy reported in the results is the
arithmetic mean of $\text{entropy}_{p,c}$ over all prompts for which at least two iterations
exist.

\subsection{Statistical Methods}

The primary hypothesis test for RQ2 is the Wilcoxon signed-rank test~\citep{Wilcoxon1945},
a non-parametric paired test that does not assume normality of the difference distribution.
This choice is appropriate given the non-normal distribution of token counts, the presence
of outliers (particularly in the Aether condition, where some prompts produce very short
contract-only outputs), and the relatively small sample size of $N = 30$ paired observations
for Gemini. The test is conducted at the $\alpha = 0.05$ significance level. For Qwen, the
Wilcoxon test is applied to prompt-level means (i.e.\ the mean \texttt{tokens\_out} across
all ten iterations for each prompt), ensuring that the unit of analysis is the prompt rather
than the individual iteration, consistent with the experimental design.

Effect size is estimated using Cohen's $d$~\citep{Cohen1988}, computed with pooled standard
deviation:
\begin{equation}
  d = \frac{\bar{x}_{\text{Vanilla}} - \bar{x}_{\text{Aether}}}{s_{\text{pooled}}}
  \label{eq:cohens_d}
\end{equation}
where $s_{\text{pooled}} = \sqrt{\frac{(n_1 - 1)s_1^2 + (n_2 - 1)s_2^2}{n_1 + n_2 - 2}}$.
Following Cohen's conventions, $|d| \geq 0.2$ indicates a small effect, $|d| \geq 0.5$
a medium effect, and $|d| \geq 0.8$ a large effect.

\subsection{Threats to Validity}

\subsubsection{Construct Validity}

The G-Suite rule engine is a custom artefact designed and implemented by the authors, which
introduces a potential bias risk if rules are defined in ways that trivially favour one
condition. This threat is mitigated by explicitly mapping each rule to a published,
peer-reviewed or industry-standard architectural specification: G1 to the Repository Pattern
(Fowler~\citep{Fowler2002PEAA}), G2 to the OWASP Secure Coding Practices~\citep{OWASP2023Secure},
and G3 to the Dependency Rule of Clean Architecture (Martin~\citep{Martin2017CleanArch}).
The weighted cost metric, while a simplification of actual billing structures, reflects the
pricing asymmetry documented across major commercial LLM API providers.

\subsubsection{Internal Validity (Cache Independence)}

Ollama's key-value (KV) cache may preserve portions of the prompt context between successive
iterations of the same prompt, potentially reducing the statistical independence of
iterations 2 through 10 for Qwen. This would mean the effective sample size for the
Qwen entropy and Wilcoxon analyses is lower than the nominal $n = 10$ repetitions per
prompt. This threat is acknowledged and partially mitigated by reporting prompt-level
means in the Wilcoxon analysis rather than treating individual iterations as independent
observations. Future work should implement explicit KV cache flushing between iterations.

\subsubsection{Internal Validity (Context Overflow Exclusions)}

Prompt P30 (implementing a \texttt{DELETE /files/:id} route) produced a context overflow
under the Qwen Vanilla condition: the first Vanilla-condition iteration exhausted the
4,096-token context window, generating a truncated, malformed response at exactly 4,096
tokens, which halted further data collection for that prompt. Only one Vanilla-condition
iteration for P30 was recorded. The single Vanilla-condition iteration for P30 is excluded
from the aggregate statistical analysis and reported separately as a qualitative finding
in \cref{sec:results}. This exclusion is conservative: it removes a data point that would
have inflated the Qwen Vanilla mean, potentially making the Aether condition appear
relatively more favourable. The exclusion is therefore unlikely to introduce a directional
bias in favour of Aether.

A single iteration of prompt P05 under the Aether condition for Llama 3.1 8B also triggered
context overflow at the 8,192-token context limit and was excluded from aggregate statistical
computation. This isolated event differs from the systematic P30 overflow observed under the
Qwen 2.5 Coder 7B Vanilla condition; with nine valid iterations retained, prompt-level mean
computation for P05 remains statistically meaningful. Corrected per-model exclusion rules
were applied throughout: P30 Vanilla was excluded for Qwen 2.5 Coder 7B only; P05 Aether
iteration 2 was excluded for Llama 3.1 8B only. No overflow events were detected for Gemma
3 12B or Qwen 2.5 Coder 14B, and no exclusions were applied to those models. By excluding only 
the single failed iteration rather than the entire prompt, the valid paired prompt count for 
Llama 3.1 8B was preserved at N=30. Gemma 3 12B and Qwen 2.5 
Coder 14B independently maintained N=30, as they experienced no overflow events.

\subsubsection{External Validity}

The study reports results for five model families spanning a parameter range from 7B to
over 100B. Although this covers both small quantized open-weight and large closed-weight
models, the generalisability of findings to other model families --- including GPT-4o,
Claude 3.5 Sonnet, Mixtral, and Code Llama variants --- cannot be assumed. The diversity
across weight types and parameter scales partially mitigates this concern. Additionally,
the 30-prompt dataset, while structurally representative of enterprise backend code
generation tasks, may not capture the full diversity of enterprise software domains.

%% =========================================================
\section{Results}
\label{sec:results}
%% =========================================================

\subsection{RQ1: Structural Correctness}

The overall structural correctness pass rate was 95.1\% across all 2,451 valid evaluations. Notably, all 121 structural failures were due to G5 (TypeScript compilation) errors occurring exclusively under the Aether condition (0\% failure in Vanilla). Specifically, Gemma 3 12B Aether achieved a 66.3\% pass rate, and Llama 3.1 8B Aether achieved 93.3\%, while all other model and condition combinations achieved 100\% correctness. Furthermore, every evaluation that produced a parseable response was
resolved on the first generation attempt; no evaluation required a second or third loop
iteration under either condition. This finding establishes an important boundary condition
for the study: the Aether Framework does not improve structural correctness relative to the
Vanilla baseline, because all five evaluated models already satisfy the
four measured architectural rules without architectural guidance, and none require
the self-healing regeneration loop to produce conforming code.

\begin{tcolorbox}[
  colback=gray!10,
  colframe=black!60,
  title=\textbf{RQ1 Summary},
  fonttitle=\bfseries,
  breakable
]
\textit{The overall structural pass rate is 95.1\% across all 2,451 valid evaluations. G5 (TypeScript compilation) failures occur exclusively under the Aether condition (0\% in Vanilla), affecting Gemma 3 12B (66.3\%) and Llama 3.1 8B (93.3\%). All other model/condition combinations achieve 100\%. Architectural pre-conditioning introduces a reliability trade-off for mid-capacity models.}
\end{tcolorbox}
%% INSERT Table 3: Summary Statistics (publication_tables.tex → Table 3)
%% INSERT Table 4: Statistical Tests / Cohen's d (publication_tables.tex → Table 4)

% --- Table 3: Summary Statistics (Mean ± SD) ---
\begin{table*}[htbp]
\centering
\caption{Summary performance metrics (Mean \(\pm\) SD) for each model under the Vanilla baseline and Aether framework conditions, ordered by capacity descending. Cells representing global minimum weighted cost are highlighted in green, and global maximum in red. All values wrapped in math mode to prevent compilation warnings.}
\label{tab:summary_stats}
\resizebox{\linewidth}{!}{
\begin{tabular}{llllll}
\toprule
\textbf{Model} & \textbf{Condition} & \textbf{Output Tokens} & \textbf{Latency (ms)} & \textbf{Total Tokens} & \textbf{Weighted Cost} \\
\midrule
Gemini 2.5 Flash & Vanilla & \(1830.27 \pm 638.42\) & \(16282.60 \pm 3998.92\) & \(2011.47 \pm 646.36\)* & \cellcolor{red!20}\(\$7502.27 \pm 2556.85\) \\
Gemini 2.5 Flash & Aether & \(872.40 \pm 863.61\) & \(25334.90 \pm 13876.22\) & \(2068.30 \pm 872.31\)* & \(\$4685.50 \pm 3453.53\) \\
Qwen 2.5 Coder 14B & Vanilla & \(749.41 \pm 138.65\) & \(58730.73 \pm 11363.41\) & \(887.45 \pm 146.46\)* & \(\$3135.69 \pm 557.76\) \\
Qwen 2.5 Coder 14B & Aether & \(494.42 \pm 235.49\) & \(42203.25 \pm 20965.54\) & \(1559.92 \pm 244.21\)* & \(\$3043.17 \pm 942.47\) \\
Gemma 3 12B & Vanilla & \(1690.85 \pm 260.21\) & \(123245.72 \pm 19576.94\) & \(1815.05 \pm 268.03\)* & \(\$6887.60 \pm 1042.40\) \\
Gemma 3 12B & Aether & \(1065.46 \pm 493.97\) & \(80702.61 \pm 37300.45\) & \(2204.36 \pm 502.54\)* & \(\$5400.74 \pm 1977.38\) \\
Llama 3.1 8B & Vanilla & \(607.66 \pm 182.39\) & \(15702.61 \pm 4602.30\) & \(726.32 \pm 189.89\)* & \(\$2549.29 \pm 731.11\) \\
Llama 3.1 8B & Aether & \(608.65 \pm 152.33\) & \(16061.84 \pm 3853.73\) & \(1650.66 \pm 160.75\)* & \(\$3476.62 \pm 612.08\) \\
Qwen 2.5 Coder 7B & Vanilla & \(448.40 \pm 160.52\) & \(20120.51 \pm 8369.05\) & \(586.64 \pm 168.39\)* & \cellcolor{green!15} \(\$1931.84 \pm 646.42\) \\
Qwen 2.5 Coder 7B & Aether & \(585.30 \pm 510.73\) & \(25299.21 \pm 23822.83\) & \(1650.92 \pm 519.58\)* & \(\$3406.83 \pm 2041.63\) \\
\bottomrule
\end{tabular}
}
\begin{flushleft}
\footnotesize{*Total tokens standard deviation is approximated. Llama 3.1 8B and Qwen 2.5 Coder 7B values reflect exclusion-corrected means (see Methodology, Context Overflow Exclusions); all other models report unadjusted raw means as no exclusions apply.}
\end{flushleft}
\end{table*}

% --- Table 4: Statistical Test Results ---
\begin{table*}[htbp]
\centering
\caption{Hypothesis testing and effect size measures comparing prompt-level average output token length under Aether vs. Vanilla, ordered descending by model parameters. Significant models (\(p < 0.05\)) are highlighted in green, non-significant in grey.}
\label{tab:statistics}
\resizebox{\linewidth}{!}{
\begin{tabular}{lllllllll}
\toprule
\textbf{Model} & \textbf{Params} & \textbf{N Pairs} & \textbf{W Stat} & \textbf{$p$-value} & \textbf{Significant} & \textbf{Cohen's $d$} & \textbf{Effect Size} & \textbf{Direction} \\
\midrule
Gemini 2.5 Flash & 100B+ & 30 & 59.0 & 1.528710e-04 & \cellcolor{green!15}Yes & 1.2613 & Large & \textcolor{green}{Aether reduces $\blacktriangledown$} \\
Qwen 2.5 Coder 14B & 14B & 30 & 56.0 & 1.105815e-04 & \cellcolor{green!15}Yes & 1.3252 & Large & \textcolor{green}{Aether reduces $\blacktriangledown$} \\
Gemma 3 12B & 12B & 30 & 7.0 & 3.539026e-08 & \cellcolor{green!15}Yes & 1.6197 & Large & \textcolor{green}{Aether reduces $\blacktriangledown$} \\
Llama 3.1 8B & 8B & 30 & 199.0 & 5.027610e-01 & \cellcolor{black!10}No & -0.0096 & Negligible & \textcolor{gray}{Negligible / No effect $\approx$} \\
Qwen 2.5 Coder 7B & 7B & 29 & 177.0 & 3.926539e-01 & \cellcolor{black!10}No & -0.3623 & Small & \textcolor{red}{Aether increases $\blacktriangle$} \\
\bottomrule
\end{tabular}
}
\end{table*}


\subsection{RQ2: Token Reduction and the Capacity Threshold}

The effect of Aether on output token generation diverges markedly across the model
families, revealing a distinct capacity-dependent transition (see \cref{tab:summary_stats} and Figure~\ref{fig:boxplot}).

For the three largest models in the study (Gemini 2.5 Flash, Qwen 2.5 Coder 14B, and Gemma 3 12B),
the Wilcoxon signed-rank test reveals a statistically significant difference between Vanilla
and Aether output token distributions (all $p < 0.001$, see \cref{tab:statistics}). The positive direction
of Cohen's $d$ (ranging from $+1.26$ to $+1.62$) indicates that the Vanilla condition
produces substantially more output tokens than the Aether condition. We classify this as the
\emph{funnel-active regime}. For these models, Aether successfully acts as a token funnel,
stripping away boilerplate and reducing output verbosity by 34--52\% (Gemini $-52.3\%$, Gemma 3 12B $-37.0\%$, Qwen 2.5 Coder 14B $-34.0\%$). 

Conversely, for the two smaller models, the effect degrades toward zero and below.
Qwen 2.5 Coder 7B shows a small negative effect ($d = -0.36$, $W=177.0$, $p=0.393$,
$N=29$), meaning Aether actually increases output token volume by 30.5\%. Llama 3.1 8B,
positioned closest to the capacity threshold, exhibits a negligible effect indistinguishable
from zero ($d = -0.010$, $W=199.0$, $p=0.503$, $N=30$), consistent with a model
operating at the lower boundary of the capacity threshold region. We classify both as
the \emph{overhead-dominant regime}.

% TABLES ARE PROVIDED IN publication_tables.tex

\begin{tcolorbox}[
  colback=gray!10,
  colframe=black!60,
  title=\textbf{RQ2 Summary},
  fonttitle=\bfseries,
  breakable
]
\textit{Architectural pre-conditioning produces a statistically significant, large-magnitude
reduction in output tokens for models at or above 12B parameters ($p < 0.001$, $d > 1.25$). 
For models at or below 8B parameters, the effect is non-significant and occasionally reverses 
direction, resulting in increased verbosity. The effect of architectural pre-conditioning on 
token generation is strongly capacity-dependent.}
\end{tcolorbox}

\subsection{RQ3: Weighted API Cost}

The weighted cost metric accounts for both the additional input token
overhead introduced by the Aether contract injection and the change in output token volume.
The arithmetic decomposition of the cost savings reveals how the funnel effect translates
into financial efficiency.

For Gemini 2.5 Flash, the mean output reduction of 958 tokens yields an output
cost reduction of $958 \times 4 = 3{,}832$ weighted units. The net weighted cost change
(accounting for the $\approx 1{,}014$ input token overhead of the contract) is therefore 
$-2{,}817$ weighted units, a 37.5\% reduction in mean weighted cost. Gemma 3 12B exhibits 
a similar dynamic, achieving a 21.6\% reduction in weighted cost.

Fascinatingly, Qwen 2.5 Coder 14B achieves near absolute cost parity ($-2.9\%$ change in weighted cost). 
For this model, the massive input token overhead of the architectural contract is almost perfectly 
offset by the reduction in output tokens. By contrast, the smaller models (Llama 3.1 8B and Qwen 2.5 Coder 7B) 
suffer massive cost penalties under Aether ($+36.4\%$ and $+76.4\%$ respectively) because the input 
overhead compounds an increase (or negligible decrease) in output generation volume.

\begin{tcolorbox}[
  colback=gray!10,
  colframe=black!60,
  title=\textbf{RQ3 Summary},
  fonttitle=\bfseries,
  breakable
]
\textit{Aether reduces weighted API cost by 21.6\% to 37.5\% for models larger than 12B parameters, 
achieves near parity ($-2.9\%$) for Qwen 14B, and massively increases costs (36.4\%--76.4\%) for models 
under 8B parameters. Cost efficiency is strictly capacity-dependent.}
\end{tcolorbox}
%% INSERT Table 5: Qualitative Code Comparison (publication_tables.tex → Table 5)

% --- Table 5: Qualitative Code Example ---
\begin{table*}[htbp]
\centering
\caption{Side-by-side comparison of generated implementations for prompt P30 (DELETE /files/:id), demonstrating qualitative architectural drift patterns in Vanilla vs. target conformance in Aether.}
\label{tab:code_example}
\begin{tabular}{p{0.48\textwidth} | p{0.48\textwidth}}
\toprule
\textbf{Vanilla Generation --- P30 (DELETE /files/:id)} & \textbf{Aether Generation --- P30 (DELETE /files/:id)} \\
\midrule
\begin{minipage}[t]{0.48\textwidth}
\scriptsize
\begin{verbatim}
// Vanilla -- Architectural Drift
import { Pool } from 'pg';
// <- raw driver (drift pattern)
import express from 'express';

const db = new Pool({
  // <- direct inst. (drift pattern)
  connectionString: process.env.DB
});

// Mock filesystem structure
// <- boilerplate (drift pattern)
const mockFiles: Record<string, {
  id: string; ownerId: string;
}> = {};

export async function deleteFile(
  req: express.Request,
  res: express.Response
) {
  // No auth check present
  // <- missing guard (drift pattern)
  const { id } = req.params;
  const result = await db.query(
    // <- raw SQL (drift pattern)
    'DELETE FROM files WHERE id=$1',
    [id]
  );
  // ... 600+ more tokens
}
\end{verbatim}
\end{minipage}
&
\begin{minipage}[t]{0.48\textwidth}
\scriptsize
\begin{verbatim}
// AETHER -- Conformant Output
import { requireUser } // Auth import
  from '../lib/auth';
import { fileRepository } // G1 Pass
  from '../lib/files';

export async function deleteFile(
  req: express.Request,
  res: express.Response
) {
  const user = requireUser(req);
  // G2 Pass

  await fileRepository // G1 Pass
    .deleteById(req.params.id, user.id);

  res.status(204).send(); // Clean
}
// 214 tokens total
\end{verbatim}
\end{minipage} \\
\bottomrule
\end{tabular}
\begin{flushleft}
\footnotesize{Note: Inline annotations indicate qualitative architectural drift patterns. Total evaluated samples N=2451.}
\end{flushleft}
\end{table*}
\clearpage

\subsection{RQ4: Generation Entropy}

Generation entropy, operationalised as the coefficient of variation across ten iterations
per prompt, reveals a model-dependent pattern across the four local open-weight models.
Under the Vanilla condition, the mean generation entropy is generally low (ranging from 0.0251
for Qwen 7B to 0.0630 for Llama 8B), indicating that models produce output of highly consistent
length when given simple, unconstrained prompts.

Under the Aether condition, entropy increases for three of the four local models. Most notably,
for Qwen 2.5 Coder 7B, entropy triples from 0.0251 to 0.0743, and for Gemma 3 12B, it more
than doubles from 0.0538 to 0.1257. Qwen 2.5 Coder 14B also shows a moderate increase
(0.0354 $\to$ 0.1065). Llama 3.1 8B is the sole exception: its entropy decreases from 0.0630
to 0.0398 under Aether --- a reduction consistent with its position at the boundary of the
capacity threshold, where the contract may exert a mild focusing effect rather than an
ambiguating one.

\begin{tcolorbox}[
  colback=gray!10,
  colframe=black!60,
  title=\textbf{RQ4 Summary},
  fonttitle=\bfseries,
  breakable
]
\textit{Aether increases generation entropy for three of the four evaluated open-weight models
(Qwen 2.5 Coder 7B: CV $0.0251 \to 0.0743$; Gemma 3 12B: $0.0538 \to 0.1257$; Qwen 14B:
$0.0354 \to 0.1065$). Llama 3.1 8B is the exception, with entropy decreasing from $0.0630$
to $0.0398$, consistent with its transitional position near the capacity threshold. The
general pattern contradicts the intuition that structural constraints should produce more
deterministic outputs.}
\end{tcolorbox}

\subsection{Qualitative Finding: P30 Context Overflow Prevention}

\begin{tcolorbox}[
  colback=blue!5,
  colframe=blue!40!black,
  title=\textbf{Qualitative Finding: P30 (DELETE /files/:id)},
  fonttitle=\bfseries,
  breakable
]
Prompt P30, which specifies the implementation of a secure file deletion route
(\texttt{DELETE /files/:id}), produced a qualitatively extreme failure under the Qwen 2.5
Coder 7B Vanilla condition. The first Vanilla-condition iteration for P30 immediately
exhausted the model's 4,096-token context window, producing a truncated and syntactically
malformed response at exactly 4,096 tokens. The overflow halted further collection: only
one Vanilla iteration was recorded for P30. The context overflow was triggered by the model's
tendency, under unconstrained generation, to emit scaffolding code, mock filesystem data
structures, and server configuration boilerplate before addressing the requested route logic.

Under the Aether condition, the same prompt was resolved correctly in 214 output tokens
across all ten iterations, representing a reduction of more than 95\% relative to the
overflowed Vanilla output. The architectural contract explicitly scoped the generation task
to a single handler function within the file management module boundary, preventing the
runaway generation behaviour observed in the Vanilla condition.

This finding represents a concrete qualitative demonstration of what we term
\emph{catastrophic output expansion prevention}: the Aether contract, even in a model
family where the aggregate statistical effect is non-significant, can prevent individual
catastrophic failures that would render generated code entirely unusable. The single P30
Vanilla iteration is excluded from all aggregate statistical analyses and reported
exclusively here.
\end{tcolorbox}

%% INSERT Figure 3: Vanilla vs. Aether Prompting Comparison (P30 Case)
\begin{figure}[H]
  \centering
  % INSERT fig3_p30_comparison below (export from architecture_diagrams.md Figure 3)
   \includegraphics[width=\linewidth]{figures/fig3_p30_comparison.pdf}
  \caption{Qualitative comparison of Vanilla (Standard Prompting) versus Aether
  (Contract-Injected) generation for prompt P30 (\texttt{DELETE /files/:id}). Under the
  Vanilla condition, the model emits raw database driver imports, direct pool instantiation,
  mock filesystem structures, and server bootstrap boilerplate, producing 4,096 tokens
  (context overflow). Under the Aether condition, the contract constrains the output to a
  single handler function with correct repository layer imports and \texttt{requireUser()}
  enforcement, resolving in 214 tokens. Illustrated output represents qualitative generation
  patterns; formal G-Suite evaluation recorded a 95.1\% overall structural pass rate.}
  \label{fig:p30}
\end{figure}

%% =========================================================
\section{Discussion}
\label{sec:discussion}
%% =========================================================

\subsection{The Funnel Hypothesis and Capacity Threshold}

The central finding of this study---that Aether produces a large, significant reduction
in output tokens for models at or above 12B parameters but a non-significant, sometimes negative
effect for models at or below 8B parameters---invites a mechanistic explanation. We propose the
\emph{funnel hypothesis}: architectural pre-conditioning acts as an output funnel only for models
that possess sufficient parametric capacity to interpret, internalise, and apply a dense structural
contract. 

A large, instruction-tuned model (e.g.\ Gemini 2.5 Flash, Qwen 2.5 Coder 14B, Gemma 3 12B) is capable
of holding the complex constraints of the Aether contract in its working memory and using them to
ruthlessly prune the generation space, stripping away unnecessary scaffolding and boilerplate. By contrast,
smaller models cross an inferred \emph{Capacity Threshold Region} between 8B and 12B parameters.
Below this threshold, the effect degrades monotonically: Qwen 2.5 Coder 7B exhibits a small
negative effect ($d = -0.36$), while Llama 3.1 8B, positioned closest to the threshold, shows
an effect indistinguishable from zero ($d = -0.010$, $p = 0.503$). This near-zero result for
Llama 3.1 8B is consistent with a model operating at the transitional boundary between the
overhead-dominant and funnel-active regimes. Below this boundary, the dense contract does not
function as a constraining filter but rather as an \emph{ambiguating stimulus}: concepts, patterns,
and requirements in the contract are reconciled at a cost of variability and verbosity.


\subsection{The Compilation Reliability Trade-off}

While the Aether Framework successfully functions as a token funnel for models at or above 12B parameters, this efficiency is partially confounded by a reduction in type-system reliability. The 121 observed G5 (TypeScript compilation) failures occurred \emph{exclusively} under the Aether condition for Gemma 3 12B and Llama 3.1 8B; no such failures occurred under the Vanilla baseline. This suggests that the dense, highly constrained nature of the architectural contract may induce hallucinated or malformed TypeScript syntax in mid-capacity models, even while it successfully forces the model to reduce its overall token volume. The result is a reliability trade-off: strong token compression coincides with a degradation in compilation success rate.

\subsection{The Entropy Paradox}

This hypothesis is corroborated by the entropy paradox. The finding that Aether increases generation
entropy for all local open-weight models is counterintuitive from a software engineering standpoint,
where structural constraints are expected to reduce the solution space and thereby produce more consistent,
predictable outputs. The observed entropy increase suggests that the architectural contract introduces
processing demands that exceed the models' reliable capacity to synthesise predictable output.

%% =========================================================
\section{Conclusion}
\label{sec:conclusion}
%% =========================================================

This paper introduced the Aether Framework, a conformance-gated generation system for
enterprise LLM code synthesis, and reported the results of a controlled empirical study
comprising 2,451 valid evaluations across 30 enterprise backend prompts, two conditions, and
five architecturally distinct model families ranging from 7B to over 100B parameters.

In answer to the four research questions posed in \cref{sec:intro}: RQ1 establishes that
Vanilla prompting achieves 100\% structural correctness, indicating that modern LLMs already
satisfy basic architectural invariants without explicit guidance, though Aether introduces a compilation reliability trade-off for mid-capacity models (95.1\% overall pass rate). RQ2 reveals that
architectural pre-conditioning produces statistically significant, large-effect reductions
in output token volume for models $\ge$ 12B parameters ($p < 0.001$, $d > 1.25$), but produces
negligible or negative effects for models $\le$ 8B parameters. RQ3 demonstrates that this
divergence translates directly to API cost: Aether yields a 21.6\%--37.5\% weighted cost
reduction for models $\ge$ 12B, achieves near parity for Qwen 14B ($-2.9\%$), and imposes
massive cost penalties for models $\le$ 8B ($+36.4\%$ to $+76.4\%$). RQ4 shows that generation
entropy increases under Aether for open-weight models, indicating that dense architectural
pre-conditioning reduces consistency.

The central contribution of this work is the first empirical demonstration that
\emph{architectural pre-conditioning exhibits strongly capacity-dependent effects}. The funnel
hypothesis provides a theoretically grounded explanation for this phenomenon: a model must possess
sufficient parametric capacity to process and apply a dense architectural contract before it can
function as an output funnel.

For enterprise software teams, the practical implication is profound. Teams utilising large managed
models or large self-hosted models ($\ge$ 12B) can deploy structural contract injection to achieve
architectural conformance while simultaneously \emph{reducing} generation costs and latency.
Conversely, teams attempting to run smaller local models ($\le$ 8B) for privacy or cost reasons
will find that dense contract injection imposes severe token overhead and output instability.
However, enterprise teams deploying Aether on mid-capacity models must carefully monitor the compilation success rate, as the framework introduces a reliability trade-off alongside its token efficiency gains.

Several directions for future investigation follow directly from these findings. Investigating the
relationship between contract complexity and the model capacity threshold required for effective
application would inform the design of adaptive contracts that scale their density to the target
model. Fine-tuning open-weight models on Aether-compliant outputs represents a promising strategy
for closing the capability gap without requiring larger base models. Finally, integration of the
Aether rule engine into an IDE language server would enable real-time, developer-interactive
conformance checking, extending the framework's utility beyond batch generation pipelines.

%% =========================================================
%% Figure Environments (placeholders)
%% =========================================================

%% INSERT Figure 4: Capacity Curve (fig4_capacity_curve.png)
\begin{figure}[H]
  \centering
  % INSERT fig4_capacity_curve below
   \includegraphics[width=\linewidth]{figures/fig4_capacity_curve.png}
  \caption{Capacity curve mapping model parameter scale (log-scale X-axis) to Cohen's $d$
  output token effect sizes. Significant models ($p < 0.05$) are shown in green, non-significant
  in grey. The yellow shaded band marks the inferred capacity threshold region (8B--12B parameters)
  separating the overhead-dominant regime from the funnel-active regime. Llama 3.1 8B sits at
  $d = -0.010$, effectively on the zero reference line, consistent with a model at the transitional
  boundary of the threshold. No interpolating trend line is shown; with five model families the
  data points are the evidence.}
  \label{fig:capacity}
\end{figure}

%% INSERT Figure 5: Token Distribution Boxplot (fig5_tokens_boxplot.png)
\begin{figure}[H]
  \centering
  % INSERT fig5_tokens_boxplot below
  \includegraphics[width=\linewidth]{figures/fig5_tokens_boxplot.png}
  \caption{Faceted distribution of output tokens (\texttt{tokens\_out}) by condition (Vanilla vs.
  Aether) across the five evaluated models, ordered by parameter capacity descending. The split-violin
  plots show observed densities cut at actual boundaries (\texttt{cut=0}) to prevent sub-zero
  probability ranges, with inner boxplots showing medians and IQRs. The extreme outlier whisker
  for Llama 3.1 8B under Aether reaches the 8,192 token limit due to a context overflow on
  prompt P05, iteration 2 (annotated and excluded from aggregate statistics).}
  \label{fig:boxplot}
\end{figure}

%% INSERT Figure 6: Weighted Cost Bar Chart (fig6_weighted_cost_bar.png)
\begin{figure}[H]
  \centering
  % INSERT fig6_weighted_cost_bar below
   \includegraphics[width=0.85\linewidth]{figures/fig6_weighted_cost_bar.png}
  \caption{Mean weighted API cost with 95\% confidence intervals across the five evaluated models.
  Weighted cost is defined as $\text{tokens\_in} + 4 \times \text{tokens\_out}$, reflecting
  pricing asymmetry. Gemma 3 12B, Qwen 2.5 Coder 14B, and Gemini 2.5 Flash exhibit cost
  reductions under Aether. Qwen 2.5 Coder 14B achieves near cost parity ($-2.9\%$), indicating
  the input token overhead is approximately offset by output reduction.}
  \label{fig:costbar}
\end{figure}

%% INSERT Figure 7: Domain Breakdown (fig7_domain_breakdown.png)
\begin{figure}[H]
  \centering
  % INSERT fig7_domain_breakdown below
   \includegraphics[height=0.90\textheight,keepaspectratio]{figures/fig7_domain_breakdown.png}
  \caption{Mean output tokens by evaluation domain (Auth, Billing, Files, Notifications, Reports, Tasks, Teams)
  and condition across the five evaluated models. Error bars represent one standard error of
  the mean. Larger capacity models show consistent token compression across all domains under
  the Aether condition.}
  \label{fig:domain}
\end{figure}


%% =========================================================
%% Declarations of Statements
%% =========================================================

\section*{CRediT Authorship Contribution Statement}
\textbf{Raduanul Haque Maruf:} Conceptualization, Methodology, Software, Data Curation, Formal Analysis, Writing -- Original Draft Preparation. \\
\textbf{Pabel Shahrear:} Supervision, Validation, Mathematics and Statistical Analysis, Writing -- Review \& Editing.

\section*{Declaration of Competing Interest}
The authors declare that they have no known competing financial interests or personal relationships that could have appeared to influence the work reported in this paper.

\section*{Declaration of Funding}
This research did not receive any specific grant from funding agencies in the public, commercial, or not-for-profit sectors.

\section*{Originality and Exclusivity}
The authors declare that this manuscript is original work, has not been published previously, and is not currently under consideration for publication elsewhere. All authors have read and approved the final version of the manuscript.

\section*{Declaration of Generative AI and AI-assisted Technologies in the Writing Process}
During the preparation of this work, the authors used Google Gemini and Anthropic Claude to improve the language, formatting, and structural clarity of the manuscript, and used an AI coding agent (Google Antigravity) to accelerate development of the evaluation pipeline. After using these tools, the authors reviewed and edited the content as needed and take full responsibility for the content of the publication.


\section*{Data Availability}

To support reproducibility and future research, the complete source code for the Aether
Framework, alongside the full benchmark dataset and statistical analysis scripts, is
publicly available in our replication package at
\href{https://github.com/Maruf-Raduan/aether-framework}{\texttt{github.com/Maruf-Raduan/aether-framework}}.

%% =========================================================
%% Bibliography
%% =========================================================
\bibliographystyle{elsarticle-num}
\bibliography{references}

\end{document}
